\chapter{Sorting}
\label{ch:sorting}

\ccprod{} sorts whole rows of the current worksheet into the order given by one
column. It is the only command in the program that rearranges more than one
cell, and the only one with an undo of its own.

\section{How the range is chosen}
\index{sorting!range}

There is no dialogue asking what to sort. With nothing selected, the cursor
decides everything:

\begin{ccsteps}
\item \textbf{The column the cursor is in} is the column sorted on.
\item \textbf{The row the cursor is on} is the first row sorted.
\item The sort runs from there down to the last row that has anything in it,
across every column that has anything in it.
\end{ccsteps}

So to sort a table with a heading row, put the cursor on the \emph{first row of
data} --- not on the heading --- in the column you want to sort by, and the
heading stays where it is.

\section{Sorting part of a sheet}
\index{sorting!a selection}

\textbf{Select a block first and only that block moves.} \key{Ctrl+A} anchors
a selection, the arrow keys extend it, and a sort then reorders those rows
within those columns and touches nothing else on the sheet.

\textbf{The block's first column is the one sorted on}, and the rest of the
block travels with it --- select \cellref{A1} to \cellref{C20} and the rows
are ordered by column~A while B and~C stay with their own row. That is the
first column of the \emph{selection}, not of the sheet: a block from
\cellref{C1} to \cellref{E9} sorts on column~C.

That is the way to sort a table that shares a worksheet with anything else:
totals below it, notes beside it, or formulas that read it.

\begin{cccaution}
If the block is narrower than the data on its rows, \ccprod{} asks
\msg{data beside the block will not move} first. Answering \key{Y} sorts the
block alone, which slides it out of step with whatever is beside it ---
sometimes exactly what is wanted, and sometimes how a table comes apart.
Widen the selection to the full width of the table to keep the rows together.
\end{cccaution}

\begin{ccfigure}
\ccscreenscale{1.1}
\begin{tikzpicture}
\node[inner sep=0pt,anchor=north west] (sb) at (0,0) {%
\ccscreenbegin{38}{8}
  \ccfillrow{0}{x16lgrey}
  \cccolhdr{0}{0}{A}{0}\cccolhdr{0}{1}{B}{1}\cccolhdr{0}{2}{C}{0}
  \foreach \r in {1,...,7}{\ccrowhdr{\r}{\r}{x16black}{x16lgrey}}
  \ccrowhdr{2}{2}{x16black}{x16yellow}
  \cccellL{1}{0}{x16black}{Item}
  \cccellL{1}{1}{x16black}{Qty}
  \cccellL{1}{2}{x16black}{Price}
  \cccursor{2}{1}
  \cccellL{2}{0}{x16black}{Widget}\cccellR{2}{1}{x16white}{40}\cccellR{2}{2}{x16black}{2.50}
  \cccellL{3}{0}{x16black}{Bolt}\cccellR{3}{1}{x16black}{12}\cccellR{3}{2}{x16black}{.75}
  \cccellL{4}{0}{x16black}{Nut}\cccellR{4}{1}{x16black}{95}\cccellR{4}{2}{x16black}{.20}
  \cccellL{5}{0}{x16black}{Washer}\cccellR{5}{1}{x16black}{7}\cccellR{5}{2}{x16black}{.05}
\ccscreenend
};
\node[inner sep=0pt,anchor=north west] (sa) at (2.62in,0) {%
\ccscreenbegin{38}{8}
  \ccfillrow{0}{x16lgrey}
  \cccolhdr{0}{0}{A}{0}\cccolhdr{0}{1}{B}{1}\cccolhdr{0}{2}{C}{0}
  \foreach \r in {1,...,7}{\ccrowhdr{\r}{\r}{x16black}{x16lgrey}}
  \ccrowhdr{1}{1}{x16black}{x16yellow}
  \cccellL{1}{0}{x16black}{Item}
  \cccellL{1}{1}{x16black}{Qty}
  \cccellL{1}{2}{x16black}{Price}
  \cccellL{2}{0}{x16black}{Washer}\cccellR{2}{1}{x16black}{7}\cccellR{2}{2}{x16black}{.05}
  \cccellL{3}{0}{x16black}{Bolt}\cccellR{3}{1}{x16black}{12}\cccellR{3}{2}{x16black}{.75}
  \cccellL{4}{0}{x16black}{Widget}\cccellR{4}{1}{x16black}{40}\cccellR{4}{2}{x16black}{2.50}
  \cccellL{5}{0}{x16black}{Nut}\cccellR{5}{1}{x16black}{95}\cccellR{5}{2}{x16black}{.20}
\ccscreenend
};
\node[anchor=north,font=\sffamily\fontsize{7.6}{9}\selectfont,text=ccink]
  at ($(sb.south)+(0,-3pt)$) {Cursor on B2. Heading row is left alone.};
\node[anchor=north,font=\sffamily\fontsize{7.6}{9}\selectfont,text=ccink]
  at ($(sa.south)+(0,-3pt)$) {After \menupath{Data > Sort up}.};
\end{tikzpicture}
\caption{Sorting on column B, ascending}
\label{fig:sort}
\end{ccfigure}

Whole rows move. The \lit{Price} column follows its \lit{Qty}; the table stays
consistent.

\section{The commands}

\begin{cctable}{Command}{Effect}{2.2in}{2.62in}
\menupath{Data > Sort up} & Sorts into ascending order \\
\menupath{Data > Sort down} & Sorts into descending order \\
\menupath{Data > Undo sort} & Puts the rows back \\
\end{cctable}

None of the three has a keyboard shortcut.

After a sort, and after an undo, the cursor is put back on \cellref{A1}.

\section{The order values sort in}
\index{sorting!order}

Values are ranked into four groups, and the groups always come in this order:

\begin{ccprocedure}[Sort order, ascending]
\item \textbf{Numbers}, smallest first. Text that reads as a number ---
\lit{456} typed as text --- sorts here as a number, not as text.
\item \textbf{Text}, alphabetically, ignoring case, on the first thirty-two
characters.
\item \textbf{Error values}, all equal to each other.
\item \textbf{Empty cells}.
\end{ccprocedure}

Descending reverses the first three groups. It does \emph{not} move the empty
rows: blanks stay at the bottom in both directions, which is almost always what
you want and is worth knowing when it is not.

A formula sorts by the value it computed, not by the formula text.

\section{Before it sorts}
\index{sorting!warnings}

\ccprod{} checks three things and may stop or ask.

\begin{cctable}{Condition}{What happens}{1.95in}{2.87in}
Nothing below the cursor & Nothing happens, and nothing is said. \\
More than 4,096 rows & \msg{too many rows to sort} \\
Not enough memory for the index & \msg{not enough memory to sort} \\
Any formula in the range & \msg{ranges over sorted rows may be wrong}, with
\key{Y} to go ahead and anything else to stop \\
\end{cctable}

\textbf{A formula carried along by a sort keeps pointing at its own row.} A
\fml{=B4*C4} sitting on row~4 becomes \fml{=B9*C9} when that row lands on
row~9 --- so a row that multiplies its own two figures goes on multiplying
them, wherever it ends up. That is what a relative reference means, and it is
what makes sorting a table with a computed column safe.

\begin{cccaution}
A \emph{range} is the case the warning is about. \fml{=SUM(B2:B10)} names ten
rows that a sort scatters, and afterwards those ten row numbers hold a
different ten values. There is no rewriting that can fix it, because the rows
the range meant are no longer next to each other.

Keep range formulas out of the rows being sorted --- in a summary block above
the table, or on another worksheet --- and sorting will not disturb them.
\end{cccaution}

\section{Undoing a sort}
\index{sorting!undo}

\menupath{Data > Undo sort} reverses the last sort exactly, by replaying the
moves backwards.

It works \textbf{once}. It refuses --- silently --- if:

\begin{ccsteps}
\item you have already undone that sort;
\item you are on a different worksheet from the one that was sorted;
\item the used range has changed size since the sort, because rows have been
added or deleted.
\end{ccsteps}

Any file command --- new, open, save, import, export --- forgets the sort as
well. Opening a menu does not.

\begin{ccnote}
\menupath{Undo sort} and \key{F5} \menupath{Undo} are two separate mechanisms
that do not interact. \key{F5} will not undo a sort, and \menupath{Undo sort}
will not undo a cell edit.
\end{ccnote}
